\chapter{Pain during Intercourse (Dyspareunia)}

\section{Definition}
Pain during intercourse (dyspareunia) is recurrent or persistent pain associated with sexual activity. It can affect men and women at any age. Causes range from physical to psychological and often involve multiple factors.

\section{Pathophysiology}
In women, dyspareunia can result from inadequate lubrication, vulvodynia, vaginismus, endometriosis, pelvic inflammatory disease, ovarian cysts, or uterine fibroids. In men, causes include prostatitis, Peyronie disease, phimosis, or urethritis. Emotional factors such as anxiety, relationship issues, or past trauma can contribute to or exacerbate pain.

\section{Etiology}
\begin{itemize}
\item Inadequate lubrication (most common cause in women)
\item Vulvodynia or vestibulodynia
\item Vaginismus (involuntary pelvic floor muscle spasm)
\item Endometriosis
\item Pelvic inflammatory disease
\item Uterine fibroids or adenomyosis
\item Ovarian cysts
\item Pelvic organ prolapse
\item Prostatitis or chronic pelvic pain syndrome (men)
\item Peyronie disease (penile curvature)
\item Urethritis or sexually transmitted infections
\item Postsurgical changes (episiotomy, hysterectomy)
\item Radiation therapy
\item Psychological factors: anxiety, depression, relationship issues, past trauma
\end{itemize}

\section{Indications for Evaluation}
\begin{itemize}
\item Pain during or after sexual intercourse
\item Superficial pain at the vaginal opening or penis tip
\item Deep pelvic pain during thrusting
\item Burning, aching, or sharp pain
\item Pain that has persisted for more than 3 months
\item Avoidance of sexual activity due to pain
\item Associated vaginal dryness, itching, or discharge
\item Relationship distress
\end{itemize}

\section{Related Symptoms}
\begin{itemize}
\item Vaginal dryness
\item Vaginismus
\item Endometriosis
\item Pelvic inflammatory disease
\item Prostatitis
\item Sexually transmitted infections
\item Menopause (vaginal atrophy)
\item Chronic pelvic pain
\end{itemize}

\section{Self-Care/Home Management}
\begin{itemize}
\item Use water-based lubricants generously
\item Practice relaxation techniques and deep breathing
\item Communicate with your partner about comfort
\item Try different positions to reduce pain
\item Engage in extended foreplay to enhance natural lubrication
\item Empty the bladder before intercourse
\item Apply a warm pack to the pelvic area before activity
\end{itemize}

\section{Diagnostic Approach}
\begin{itemize}
\item Detailed sexual and medical history
\item Pelvic examination (women)
\item Genital examination (men)
\item Testing for sexually transmitted infections
\item Pelvic ultrasound if endometriosis or fibroids suspected
\item Laparoscopy for suspected endometriosis
\item Psychological evaluation if indicated
\end{itemize}

\section{Treatment/Management Overview}
\begin{itemize}
\item Treat underlying cause (infection, endometriosis, prostatitis)
\item Pelvic floor physical therapy
\item Topical estrogen for vaginal atrophy
\item Lidocaine or other topical analgesics
\item Cognitive behavioral therapy or sex therapy
\item Vaginal dilators for vaginismus
\item Couples counseling
\item Medication adjustments (change contraceptives or antidepressants)
\end{itemize}
